% ID: MAT-ALG-PER-005
% Titolo: Sconti e percentuale complessiva
% Disciplina: Matematica
% Area: Algebra
% Argomento: Percentuali
% Classe: Prima secondaria di secondo grado
% Difficolta: 3
% Tipo: laboratorio
% Tag: matematica, algebra, percentuali
% Fonte: laboratorio_percentuali_anonimizzato
% Provenienza: materiale fornito dal docente; testo anonimizzato; soluzione redatta per Exergo

\begin{esercizio}
In un foglio di calcolo inserisci i prezzi e gli sconti:
\[
\begin{array}{c|rr}
\text{Articolo}&\text{Prezzo (euro)}&\text{Sconto}\\\hline
\text{Quaderno}&4&10\%\\\text{Zaino}&50&20\%\\
\text{Calcolatrice}&24&15\%\\\text{Cuffie}&36&25\%\\\text{Mouse}&20&5\%
\end{array}
\]
Calcola risparmio e prezzo finale di ciascun articolo, i totali e lo sconto complessivo.
Confronta l'articolo con lo sconto percentuale maggiore e quello con il risparmio maggiore in euro.
\end{esercizio}

\begin{soluzione}
I risparmi sono $0{,}40;10;3{,}60;9;1$ euro; i prezzi finali $3{,}60;40;20{,}40;27;19$ euro.
Totali: prezzo iniziale 134 euro, risparmio 24 euro, prezzo finale 110 euro.
Lo sconto complessivo è $(24/134)100\%\simeq17{,}91\%$.
In Excel, prezzi in B2:B6 e sconti in C2:C6: D2 = \texttt{=B2*C2}, E2 = \texttt{=B2-D2}, poi copia in basso.
La percentuale complessiva è \texttt{=SOMMA(D2:D6)/SOMMA(B2:B6)}.
Le cuffie hanno lo sconto percentuale maggiore, lo zaino il risparmio in euro maggiore.
I prezzi diversi impediscono di usare la media aritmetica degli sconti.
\end{soluzione}
